\section{A return set determined by stationary data}
\label{sec:stationary-return-set}
Retain two disjoint column bases, half noise $P$, $\|\pi\|_{15/7}\le B$,
and stationary matching through ambient order $k$. The dimension is unrestricted.
Define $V=P\sqrt{\pi}$ on the quotient, or equivalently pull it back to the cube.
Fix $t=\sigma k$ as in (FF1), choose $0<a<1/153$, and put
\[
 Z=\{V<e^{-at}\}.
\]
Write $s_E$ for the top killed eigenvalue of $K$ on $E$ and $R_E$ for the
largest conditional stationary one-step return to any nonempty subset of $E$.
Both quantities are zero for the empty set.

\paragraph{Theorem.}
There are positive $c,C$ depending only on $B,\sigma,a$ such that
\[
 |s_Z-s_{\{U>e^t\}}|\le Ce^{-ck},\qquad
 s_Z/4-Ce^{-ck}\le\Delta\le s_Z+Ce^{-ck}.                 \tag{SR1}
\]
Consequently full ordered spectral recovery is equivalent to $R_Z\to0$.
The set to inspect can be computed from the stationary density and the ordinary
noise alone. Evaluating the return still uses the transformed chain.
In particular if every noise row hits $\{\pi\ge b\}$ with probability at least
$\eta$ and $\log(1/(\eta\sqrt b))=o(k)$ for $0<b,\eta\le1$, recovery is exponential.
Indeed $V\ge\eta\sqrt b$ and $Z$ is eventually empty.

\paragraph{Why the visible part of a high multiplier can be discarded.}
By Cauchy--Schwarz in each noise row,
\[
 V^2=(P\sqrt{hg})^2\le (Ph)(Pg)=g\,Pg.
\]
Here $\|g\|_2^2=\langle h,P^2h\rangle\le\langle h,Ph\rangle=1$ by the
spectrum of $P$. On quotient functions the two bases give consecutive
contractive forward steps $2\to5\to17$ by the Boolean
\href{https://doi.org/10.5802/aif.357}{hypercontractive inequality}.
Any unused columns act contractively. Thus $\|Pg\|_{17}\le1$.
On $Z^c$, $g^{-1}\le e^{2at}Pg$. With $1/w=7/15+2/17=149/255$,
H\"older yields $\|U1_{Z^c}\|_w\le B e^{4at}$ from $U=\pi/g^2$.
Put $A=\{U>e^t\}$ and $v=5/3$. Truncation gives
\[
 \|U1_{A\setminus Z}\|_v\le B^{153/149}
           \exp[-(4-612a)t/149].                            \tag{SR2}
\]
Indeed $w/v=153/149$. This estimate also gives the restricted reciprocal
bound $\|g^{-1}1_{Z^c}\|_{17}\le e^{2at}$ directly, without a bound on
the mass of edges between $Z$ and its complement.

\paragraph{Why low multipliers inside the new set can be added.}
The uniform mass of $\{\pi\ge1/2\}$ is at least $(1/(2B))^{15/8}$
by H\"older and $\mathbb E\pi=1$. For $f=\sqrt\pi$, the generalized
norm $\|f\|_{7/8}$ is therefore bounded below by a positive constant
depending only on $B$. The \href{https://arxiv.org/abs/1108.1210}{reverse
hypercontractive} steps $7/8\to1/2\to-1$ give
$\mathbb E(1/V)\le C_B$. Markov's inequality yields $\mu(Z)\le C_B e^{-at}$.
Section~\ref{sec:floor-free} proves $\|U1_{A^c}-\pi\|_{5/3}\le Ce^{-ck}$.
Since $1/v-7/15=2/15$, it follows that
\[
 \|U1_{Z\setminus A}\|_v
 \le Ce^{-ck}+B\mu(Z)^{2/15}\le C'e^{-c'k}.               \tag{SR3}
\]
For $Q=P^{1/2}$ write $S_E=QM_{U1_E}Q$. Its norm equals $s_E$ because
$LL^*$ and $L^*L$ have the same nonzero spectrum for $L=QM_{\sqrt{U1_E}}$.
The quotient weight bound of Section~\ref{sec:floor-free} gives
$\|QM_WQ\|\le\|W\|_v$ for nonnegative quotient weights.
The triangle inequality with (SR2) and (SR3) now bounds
$\|S_A-S_Z\|$ exponentially. Apply (FF1) and the killed-chain return
inequalities $R_Z\le s_Z\le\sqrt{2R_Z-R_Z^2}$, which hold for any set.
This proves (SR1), including when one set is empty.

\paragraph{Recovery across holes of superlinear column radius.}
Use the repeated identity bases and full-rank perturbation of
Section~\ref{sec:coverage}, now with $m=k^2$ blocks of length $\ell=k+1$,
$n=2^m$, $r\ge m\ell+3$, and $\varepsilon=e^{-n^2}$. On parity words
of weight at most $\lfloor m/4\rfloor$ put density $\varepsilon$, on
the all-ones word put $n^{7/15}$, and on the rest put the constant
that makes the average one. The exceptional ball occupies $o(1)$ of
the words by the binomial tail, and the spike contributes
$n^{-8/15}$ to the mean. The bulk value tends to one. Apply the same
mean-zero full-span perturbation of magnitude at most $\varepsilon/4$.
Matching follows from ambient degrees at least $2\ell$ and $2(r-1)$.
The critical moment remains bounded; every fixed larger positive moment
and every fixed inverse moment diverges by the spike and the zero word,
respectively.

Starting at any syndrome the output block parities are independent, each
with bias of absolute value $4^{-\ell}$, since every syndrome coordinate
receives two independent half-noise factors. The probability of each
parity word is therefore at least $2^{-m}(1-4^{-\ell})^m$.
Uniformly in the starting point, the probability of the bulk tends to
one. There $\pi\ge1/2$ for all large $k$, so (SR1) gives exponential
recovery for every indicated rank. Yet from the zero word a column
translation changes at most one parity. Every map of radius $o(k)$
leaves that starting syndrome in the low ball, with density at most
$5\varepsilon/4$. No threshold with $\log(1/b)=o(k)$ can include
that image. Thus the preceding section's local map hypothesis fails
at every admissible radius and threshold, regardless of allowed load.
This separation proves the advantage of using the noise distribution
over a prescribed column radius in this family. Necessity of a
uniform lower bound on $V$, and unconditional critical recovery at
arbitrary rank, do not follow.

\paragraph{A nonempty stationary return set can be harmless.}
For $r=k^3$ take two identity bases on the syndrome cube and let
$M=\lceil 2k\log_2 r+2\rceil$ for $k\ge2$. There exists a binary linear
subspace $C$ spanned by at most $M$ vectors whose orthogonal space contains
no nonzero word of weight at most $k$. Indeed choose $M$ independent uniform
vectors (linear independence is unnecessary). Each fixed nonzero word is
orthogonal to all of them with probability $2^{-M}$, and there are at most
$\sum_{i=1}^k r^i\le 2r^k$ words to exclude. Fix a choice avoiding them all.
Let $E=C+\{y:|y|\le\lfloor 2r/5\rfloor\}$. Under uniform probability,
$\mu(E)\le 2^M\exp(-c r)$ by the binomial tail and a union bound.
Since $M=o(r)$, this tends to zero exponentially in $r$. Put
$\pi=e^{-2r}$ on $E$ and $\pi=(1-e^{-2r}\mu(E))/(1-\mu(E))$ elsewhere,
for all sufficiently large $k$. The density is positive, normalized and
bounded by two. Its invariance under translations by $C$ puts all its
nonconstant Walsh support in $C^\perp$, hence its stationary moments match
through order $k$ even before doubling degrees on the ambient cube.
Positive symmetric scaling realizes it as a stationary density.

The raw syndrome row from zero has independent bits with flip probability
$(1-1/4)/2=3/8$. Its chance to leave $E$ is at most the upper binomial tail
above $2r/5$, hence at most $\exp(-c' r)$. Thus
$V(0)\le e^{-r}+\sqrt2\,\exp(-c'r)$. In particular $Z$ contains zero
for every fixed admissible $a,\sigma$ and large $k$. Nevertheless (P1)
gives exponential recovery, using for example $p=3$ and two bases.
This rules out necessity of uniform one-step visibility even with bounded
stationary densities. It leaves open the possibility of a returning subset
of $Z$ at the critical moment in general.
